\section{Mở đầu}
\label{sec:intro}

\subsection{Bối cảnh và động cơ}
Kubernetes đã trở thành nền tảng điều phối container phổ biến nhất cho điện toán
đám mây hiện đại. Trong mô hình triển khai \emph{multi-tenant} (nhiều người thuê),
các tenant được cô lập \emph{mức logic} bằng không gian tên (namespace), hạn ngạch
tài nguyên và chính sách mạng---nhưng tất cả vẫn \emph{chia sẻ chung một nhân} hệ
điều hành của nút (node). Khác với máy ảo, ranh giới cô lập giữa các tenant vì thế
mỏng hơn nhiều và hội tụ tại một điểm duy nhất: giao diện lời gọi hệ thống
(syscall)---ranh giới giữa không gian người dùng và nhân~\cite{ref1_jarkas_survey,
ref3_sultan_survey}.

Hệ quả là \emph{leo thang đặc quyền} (privilege escalation) và \emph{thoát
container} (container escape) trở thành lớp mối đe dọa nguy hiểm nhất trong môi
trường multi-tenant: một tenant bị xâm phạm có thể lạm dụng syscall đặc quyền để
chiếm quyền \texttt{root} trên nút và phá vỡ cô lập, ảnh hưởng tới mọi tenant khác
trên cùng nút~\cite{ref5_combe_docker}. Các họ lỗ hổng tiêu biểu---PwnKit
(\mbox{CVE-2021-4034}), Dirty~Pipe (\mbox{CVE-2022-0847}), thoát \texttt{runc}
(\mbox{CVE-2019-5736})---đều để lại \emph{dấu vết hành vi} ở tầng syscall (gọi
\texttt{ptrace}, \texttt{setuid}/\texttt{capset}, \texttt{unshare}/\texttt{setns},
\texttt{mount}, nạp mô-đun nhân, \texttt{bpf}\dots). Các cơ chế tĩnh như hồ sơ
seccomp khó duy trì, dễ chặn nhầm thao tác hợp lệ hoặc bỏ sót tấn công nhiều bước,
và không nhận thức được \emph{chuỗi} hành vi đặc trưng cho leo thang đặc quyền.

Một hướng tiếp cận thích nghi đã được công trình tham chiếu
DeSFAM~\cite{desfam_original} đề xuất: kết hợp lập hồ sơ syscall lai (tĩnh + động),
phát hiện bất thường bằng học máy độ trễ thấp (mô-đun \emph{SyscallAD} dùng VAE và
iForest), tính điểm rủi ro theo ngữ cảnh (ánh xạ MITRE ATT\&CK, tương quan CVE) và
thực thi trong nhân qua eBPF/LSM. Tuy nhiên DeSFAM đánh giá trên hạ tầng riêng,
không công bố mã nguồn đầy đủ và không đặt trọng tâm vào ngữ cảnh \emph{multi-tenant
Kubernetes}. Chuyên đề này kế thừa thiết kế lõi phát hiện của DeSFAM làm
\emph{tham chiếu} và xây dựng một cơ chế phát hiện leo thang đặc quyền vận hành
trực tiếp trong Kubernetes.

\subsection{Mục tiêu và phạm vi báo cáo}
Báo cáo chuyên đề này trình bày một \emph{cơ chế phát hiện tấn công leo thang đặc
quyền} cho môi trường multi-tenant Kubernetes dựa trên phân tích syscall, xây dựng
trên nền các nghiên cứu liên quan---trọng tâm là công trình tham chiếu
DeSFAM~\cite{desfam_original}. Đây là một báo cáo môn học theo hướng nghiên cứu
chủ đề: hệ thống hóa cơ sở lý thuyết, thiết kế và hiện thực hóa cơ chế, rồi đánh
giá thực nghiệm. Cụ thể, báo cáo:
\begin{enumerate}[label=(\roman*)]
  \item Xây dựng lõi phát hiện bất thường syscall (VAE + iForest, hợp nhất tập
        hợp) theo thiết kế tham chiếu SyscallAD của DeSFAM~\cite{desfam_original},
        trong đó các đặc trưng \emph{hiện diện của syscall đặc quyền} đóng vai trò
        tín hiệu chính cho leo thang đặc quyền; huấn luyện và kiểm chứng ngoại tuyến
        trên bộ dữ liệu công khai DongTing~\cite{ref26_dongtingdataset};
  \item Hiện thực hóa cơ chế thành một bộ phát hiện \emph{trực tiếp} trong
        Kubernetes: thu thập syscall theo namespace (tenant) bằng
        Tetragon~\cite{tetragon} (eBPF), trích đặc trưng theo cửa sổ trượt giống hệt
        lúc huấn luyện, chấm điểm và phát cảnh báo qua Prometheus/Grafana;
  \item Đánh giá khả năng phân tách lành tính/tấn công trên các kịch bản leo thang
        đặc quyền mô phỏng theo CVE thực, và đối chiếu trung thực với số liệu công
        bố của công trình tham chiếu.
\end{enumerate}

\subsection{Nội dung và bố cục báo cáo}
Báo cáo tập trung vào lớp phát hiện (không bao gồm khâu chặn nội tuyến trong nhân)
và trình bày các nội dung chính sau:
\begin{itemize}[leftmargin=1.6em]
  \item Một cơ chế phát hiện leo thang đặc quyền cho multi-tenant Kubernetes dựa
        trên phân tích syscall, với quy trình huấn luyện \emph{hoàn toàn bằng
        Docker}, có thể chạy lại từ đầu và lưu vết kết quả số (\texttt{results.json}).
  \item Một bộ phát hiện trực tiếp dựa trên Tetragon/eBPF giám sát theo namespace
        trên Kubernetes, cho thấy thiết kế tham chiếu DeSFAM khả thi trong môi
        trường vận hành thực.
  \item Một đối chiếu \emph{trung thực} giữa cơ chế đề xuất và số liệu công bố của
        DeSFAM, cùng bài học về căn chỉnh ``bảng chữ cái cảm biến'' (sensor alphabet)
        giữa huấn luyện và phục vụ.
\end{itemize}

Phần còn lại của báo cáo được tổ chức như sau: Mục~\ref{sec:related} trình bày cơ
sở lý thuyết và công trình liên quan; Mục~\ref{sec:methodology} mô tả thiết kế cơ
chế phát hiện; Mục~\ref{sec:setup} trình bày triển khai và thiết lập thực nghiệm
bằng Docker/Kubernetes; Mục~\ref{sec:evaluation} báo cáo kết quả và đánh giá (bao
gồm thảo luận và hạn chế); Mục~\ref{sec:conclusion} kết luận.
